\documentclass[11pt,a4paper]{article}
\usepackage[margin=1in]{geometry}
\usepackage[T1]{fontenc}
\usepackage[utf8]{inputenc}
\usepackage{lmodern}
\usepackage{microtype}
\usepackage{enumitem}
\usepackage{xcolor}
\usepackage{hyperref}
\hypersetup{colorlinks=true, linkcolor=black, urlcolor=blue}
\setlength{\parindent}{0pt}
\setlength{\parskip}{0.55em}
\setlist[itemize]{leftmargin=1.5em,itemsep=0.2em,topsep=0.2em}
\setlist[enumerate]{leftmargin=1.7em,itemsep=0.2em,topsep=0.2em}
\pagenumbering{arabic}
\begin{document}
\begin{center}
{\LARGE\bfseries Proposta para Painel Web de Gestão Operacional para Despachantes Veiculares}
\end{center}
\vspace{0.5em}

Olá,\par

Meu nome é Fabricio Santana. Sou engenheiro de software e líder técnico com mais de 20 anos de experiência em desenvolvimento backend, arquitetura de software, integração com bancos de dados, testes automatizados, computação em nuvem e entrega de sistemas críticos. Também liderei equipes grandes e atuei em projetos em que consistência de dados, regras operacionais, segurança, validação e manutenção são fatores centrais.\par

Seu projeto tem forte aderência ao meu perfil porque o desafio principal não é apenas construir uma interface moderna. A parte importante é transformar os dados existentes do SDG em uma Torre de Controle operacional que realmente ajude a equipe a acompanhar prioridades, gargalos, processos parados e produtividade sem virar um novo sistema transacional.\par

Minha abordagem para esse tipo de projeto combina arquitetura web, integração com fonte de dados existente, definição de indicadores, controle de acesso, atualização automática, experiência responsiva e handoff técnico em um plano de entrega coerente. Isso é importante aqui porque a qualidade do resultado depende de definir bem a fonte de verdade dos dados, as regras dos KPIs e o limite do primeiro release operacional.\par

\section*{Metodologia de trabalho}

Meu método de entrega foi desenhado para reduzir risco antes do desenvolvimento principal. O trabalho começa com uma Discovery Sprint fechada e depois segue para um pacote de entrega somente quando escopo, prioridades, dependências, responsabilidades e critérios de aceite estiverem claros.\par

O primeiro passo recomendado é uma Discovery Sprint de 1 semana, com valor fixo de R\$ 3.000. Nessa fase, posso conduzir até 5 reuniões remotas para entender o fluxo operacional do despachante, revisar como o SDG expõe os dados, mapear indicadores, riscos, dependências, perfis de acesso, contexto de exibição e limite do primeiro release operacional. O entregável é um plano prático de implementação com escopo recomendado, abordagem técnica, marcos, critérios de sucesso, riscos, responsabilidades e recomendação do pacote de entrega.\par

Depois da discovery, o projeto pode seguir para um dos pacotes fixos de 4 semanas. O pacote é escolhido conforme complexidade, risco, qualidade da fonte de dados, necessidade de consolidação, exigências de visualização e cuidados de segurança e observabilidade.\par

\textbf{Opções de pacote}\par

\begin{itemize}
\item \textbf{Essential Delivery} --- R\$ 12.000, 4 semanas. Adequado apenas se a discovery mostrar um recorte bem estreito, com poucos indicadores e fonte de dados muito simples.
\item \textbf{Product Acceleration} --- R\$ 24.000, 4 semanas. Adequado para um painel operacional com dashboard principal, prioridades, gargalos, processos parados, produtividade, fluxo visual, login e atualização automática sobre uma fonte de dados existente.
\item \textbf{Scale \& Intelligence} --- R\$ 36.000, 4 semanas. Adequado quando a integração exigir maior consolidação, tratamento de qualidade de dados, observabilidade reforçada, múltiplas fontes ou maior sofisticação operacional.
\end{itemize}

Para esta oportunidade, a recomendação inicial mais provável é \textbf{Product Acceleration}, com possibilidade de revisão para \textbf{Scale \& Intelligence} se a discovery mostrar maior complexidade de integração ou dados.\par

\textbf{Garantia de defeitos de construção.} Após entrega e aceite, incluo garantia de 30 dias para defeitos de construção relacionados ao escopo aprovado. Isso cobre erros de implementação, comportamentos combinados que deixem de funcionar e defeitos introduzidos pelo código entregue. Não cobre novas funcionalidades, mudanças de regra operacional, alterações na estrutura ou disponibilidade da fonte de dados do SDG, mudanças em serviços de terceiros, infraestrutura fora do projeto ou solicitações que ampliem os critérios de aceite aprovados.\par

\section*{Análise da oportunidade}

Este projeto deve ser tratado como um sistema de acompanhamento gerencial sobre dados existentes, e não como um novo ERP ou CRM. O painel precisa ajudar a equipe a enxergar o que precisa de atenção agora, onde o fluxo está travando, quais processos estão parados e como a produtividade está distribuída, sem permitir qualquer cadastro, alteração ou movimentação operacional.\par

Os módulos pedidos são coerentes, mas o risco real está em como os dados serão lidos e interpretados:

\begin{itemize}
\item o dashboard principal depende de definições claras para atraso, vencimento, conclusão e falta de movimentação;
\item o painel de prioridades depende de regras objetivas para o que merece atenção imediata;
\item os gargalos operacionais dependem de um fluxo de etapas representado de forma consistente na fonte;
\item a produtividade da equipe depende de associação confiável entre processos, atividades e responsáveis;
\item a exibição em TV corporativa pede hierarquia visual, refresh controlado e leitura fácil à distância;
\item a camada de leitura precisa ser estritamente segregada da operação transacional.
\end{itemize}

Também é importante tratar a integração com o SDG como risco central desde o início. O sucesso do projeto depende de confirmar se a fonte de dados virá por API, leitura de banco ou exportação automática, qual é a frequência de atualização viável, como lidar com atrasos de sincronização e que nível de qualidade e completude os dados realmente oferecem.\par

Outro ponto importante é a definição dos indicadores. Um painel desse tipo só funciona bem quando atraso, processo parado, prioridade, etapa atual e tempo médio possuem regras alinhadas com a operação real do cliente. Sem isso, o sistema pode ficar bonito, mas pouco útil para decisão diária.\par

Por fim, a exigência de manter o painel somente leitura precisa virar requisito técnico explícito. A arquitetura deve impedir que a solução seja usada para editar cadastros, emitir documentos ou movimentar processos, preservando seu papel de Torre de Controle e reduzindo risco operacional.\par

\section*{Plano de trabalho proposto}

Eu dividiria o trabalho nas seguintes fases:\par

\textbf{1. Discovery e desenho funcional}\par

\begin{itemize}
\item Revisar como o SDG expõe os dados e validar a melhor estratégia de leitura.
\item Mapear o fluxo operacional real do despachante e as etapas que devem aparecer no painel.
\item Definir os KPIs, regras de prioridade, critérios de atraso e conceito de processo parado.
\item Delimitar o primeiro release operacional e os contextos de exibição: desktop, mobile, tablet e TV.
\item Produzir um plano de implementação com arquitetura, escopo, riscos, critérios de aceite e pacote recomendado.
\end{itemize}

\textbf{2. Arquitetura de leitura e base técnica}\par

\begin{itemize}
\item Definir a arquitetura web, autenticação, perfis de acesso e estratégia de deploy.
\item Estruturar a camada de ingestão ou sincronização de dados em modo somente leitura.
\item Definir tratamento para falhas de carga, inconsistências e atrasos de atualização.
\item Preparar a base técnica para responsividade, auto refresh e exibição em TV corporativa.
\end{itemize}

\textbf{3. Dashboard principal e prioridades}\par

\begin{itemize}
\item Implementar os indicadores centrais de processos em andamento, concluídos, atrasados, próximos do vencimento e sem movimentação.
\item Implementar o painel de prioridades com foco no que exige atenção imediata.
\item Validar filtros, contadores e regras de negócio com base nos dados reais.
\end{itemize}

\textbf{4. Gargalos, processos parados e produtividade}\par

\begin{itemize}
\item Implementar a visão de gargalos operacionais por etapa.
\item Implementar a lista de processos parados com número, etapa, última movimentação, dias sem atualização e responsável.
\item Implementar a produtividade da equipe com métricas auditáveis por colaborador.
\item Entregar o fluxo operacional visual com quantidade por etapa, tempo médio de permanência e atrasos.
\end{itemize}

\textbf{5. Testes, publicação, documentação e handoff}\par

\begin{itemize}
\item Testar fluxos principais, login, atualização automática, comportamentos de leitura e consistência das métricas.
\item Publicar a aplicação em ambiente acessível por computador, celular, tablet e TV corporativa.
\item Documentar arquitetura, integração de dados, limites conhecidos e operação básica.
\item Fazer handoff técnico e orientação inicial para uso seguro da ferramenta.
\end{itemize}

\section*{Prazo estimado}

Minha recomendação imediata é 1 semana de Discovery Sprint. Depois disso, o primeiro ciclo de implementação pode seguir em 4 semanas, focado no primeiro release operacional acordado.\par

Cronograma prático sugerido:\par

\begin{itemize}
\item Semana 1: Discovery Sprint, validação da fonte de dados, definição dos KPIs, regras de prioridade e plano de implementação.
\item Semanas 2 a 5: primeiro ciclo de implementação e entrega do painel operacional inicial.
\end{itemize}

Se a discovery mostrar maior complexidade de integração, baixa qualidade dos dados ou necessidade de consolidação mais pesada, a recomendação pode passar para um pacote maior ou para mais de um ciclo de implementação.\par

\section*{Disponibilidade e comunicação}

Posso atuar remotamente, com comunicação em português, por mensagens, chamadas agendadas e atualizações escritas. Prefiro checkpoints curtos ao longo da semana, especialmente após a validação técnica da fonte de dados, a definição dos KPIs e os principais marcos de implementação.\par

Isso ajuda a manter o projeto alinhado com o fluxo real da operação, com as prioridades do time e com a confiabilidade dos indicadores exibidos.\par

\section*{Orçamento}

Com base no entendimento atual, meu próximo passo recomendado é:\par

\begin{itemize}
\item \textbf{Discovery Sprint} --- R\$ 3.000, 1 semana.
\end{itemize}

Depois da discovery, o pacote mais provável para o primeiro release é:\par

\begin{itemize}
\item \textbf{Product Acceleration} --- R\$ 24.000, 4 semanas, caso o primeiro release inclua dashboard principal, prioridades, gargalos, processos parados, produtividade, fluxo visual, login e atualização automática sobre uma fonte de dados acessível.
\end{itemize}

Esta estimativa assume que o cliente fornecerá acesso técnico ao SDG ou a sua camada de exportação, exemplos reais de dados, definição das regras operacionais, critérios de prioridade, contexto de uso em TV corporativa, domínio ou acesso ao provedor de hospedagem e rapidez nas decisões sobre o primeiro release. Custos recorrentes de hospedagem, domínio, licenças de terceiros, infraestrutura de monitoramento e serviços externos não estão incluídos no valor de desenvolvimento.\par

Se a revisão mostrar que a integração exige saneamento mais pesado de dados, múltiplas fontes ou maior observabilidade, recomendarei um pacote superior ou ciclos adicionais em vez de prometer um escopo maior do que cabe no primeiro release.\par

\section*{Perguntas antes da confirmação final}

Antes de confirmar escopo e marcos finais, eu gostaria de alinhar:\par

\begin{enumerate}
\item O SDG da Bludata já oferece API documentada, acesso de leitura ao banco ou exportação automática estruturada?
\item Quem vai fornecer e aprovar o acesso técnico necessário a essa fonte de dados?
\item Quais campos hoje permitem calcular etapa atual, última movimentação, prazo, atraso, conclusão e responsável?
\item Como o fluxo operacional é representado hoje dentro do SDG: existem status padrão ou a equipe usa convenções próprias?
\item O que define exatamente um processo parado, um processo atrasado e um processo prioritário para o cliente?
\item A atualização precisa ser em tempo quase real, por janelas programadas ou por cargas periódicas?
\item Haverá um único perfil de visualização ou diferentes perfis de acesso?
\item A exibição em TV corporativa precisa de modo fullscreen, rotação automática ou layout específico?
\item Qual volume aproximado de processos e histórico precisa ser suportado no lançamento?
\item Já existe um orçamento aprovado e uma data-alvo para a primeira versão operacional?
\end{enumerate}

Com essas respostas, consigo fechar o plano de entrega e a recomendação final de implementação.\par

\end{document}
